Learners embedded in a population adapt not only what they predict but whom they
learn from, so a representation error in one agent can restructure the
information flow of the whole group. We study a population of neural agents
whose perceptron opinions and directed trust variables co-adapt through an
entropic approximation to Bayesian online learning derived for a noisy
student--teacher channel. A fraction $\fd$ of agents extend their representation
of an incoming message with one coordinate that carries no information about the
issue and depends only on the sender's class, shifting the agreement field that
drives their update by a strength $d$. We show that this shift translates the
local opinion--trust flow rigidly, breaking an exact $\mathbb{Z}_2$ symmetry of
the unbiased dynamics and displacing a separatrix we obtain in closed form, and
we derive a parameter-free prediction $\CcT=\fd\,\Gresp(d)$ for the trust--class
correlation. Replicated simulations exceed it by \num{residual}, which measures
the amplification the population itself contributes. Sweeping $(d,\fd)$ with
\num{24} replicates per point reveals \num{n} collective regimes, separated by
crossovers rather than transitions; magnitude-matched controls whose sign is
uncorrelated with class establish that the effect requires class correlation and
not merely added perturbation. Varying the agenda complexity $\alpha=P/K$
reverses whether affective or ideological structure forms first, crossing at
$\alpha=\num{crossover}$. We further prove that the apparent transition at $d=0$
is a forced zero-crossing of an odd response rather than a bifurcation: the
affective sector has no linear restoring force at the class-symmetric state, so
no critical field strength exists at any population size.
